\section{Conclusion and Future Work}
\label{sec:conclusion}

\subsection{Direct Answers to Research Questions}
\label{sec:concl_summary}
This investigation examined whether surgically integrating 2D selective state-space modeling into lightweight single-stage detectors provides a viable strategy for warehouse perception under edge constraints. Our empirical findings directly answer the three guiding research questions:

\begin{enumerate}
    \item \textbf{RQ1 (Architectural Efficiency vs. Accuracy Maintenance):} Substituting the $P_5$ attention block with \texttt{WarehouseSS2DBlock} eliminates $63.9\%$ of the parameters in that stage ($357{,}632$ vs. $990{,}976$) and reduces total network parameters by $6.72\%$ ($8.80\text{ M}$ vs. $9.43\text{ M}$) at $21.2\text{ GFLOPs}$. However, accuracy maintenance is nuanced: while forklift detection improves markedly ($+6.1\text{ pp}$ to $+9.6\text{ pp}\text{ mAP}_{50}$) and macro recall rises by $+7.3\text{ pp}$ ($64.5\%$ vs. $57.2\%$), precision declines by $9.3\text{ pp}$ (leaving overall F1 virtually identical at $0.683$ vs. $0.673$) and strict $\text{mAP}_{50-95}$ drops by $2.6\text{ pp}$ to $4.3\text{ pp}$. Furthermore, while 4-class global $\text{mAP}_{50}$ is lower ($67.10\%$ vs. $69.75\%$), this gap is primarily driven by the sparse pallet truck tail class ($n=13$); on primary classes (\texttt{worker}, \texttt{forklift}, \texttt{amr}), 3-class macro $\text{mAP}_{50}$ reaches $82.61\%$ to $84.64\%$, surpassing the baseline control ($82.13\%$). Thus, SS2D operates as an operating-point shift favoring safety recall over fine localization rather than an unambiguous Pareto improvement.
    \item \textbf{RQ2 (Optimization Stability and Staged Warmup):} Zero-$\gamma$ affine batch normalization residual initialization ($\gamma_0 = \mathbf{0}, \beta_0 = \mathbf{0}$) mathematically guarantees identity pass-through at step $0$, eliminating initialization disruption when inserting uncalibrated SSM weights into pretrained backbones. However, because our two-stage Arm 3 protocol coupled backbone freezing with a refined learning rate ($\eta = 0.0005$ vs. $0.00125$), its stabilization represents a compound effect of scheduled warmup and learning rate damping.
    \item \textbf{RQ3 (Edge Viability and Latency Realism):} Physical embedded deployment on a Raspberry Pi 4 Model B demonstrates that the zero-plugin ONNX model satisfies embedded memory constraints ($371.9\text{ MB}$ RSS, maintaining a $63.7\%$ margin under $1024\text{ MB}$ / $62.8\%$ under $1000\text{ MB}$). However, unrolled sequential recurrence on general-purpose CPUs yields $0.72\text{ FPS}$ ($1398.19\text{ ms}$), running $19.5\%$ slower than baseline convolution ($1169.73\text{ ms}$) due to operator dispatch overhead across $14{,}767$ nodes. This proves that CPU deployment provides memory compliance and software portability rather than line-rate throughput, establishing that real-time edge tracking ($>30\text{ FPS}$) requires parallel prefix-scan engines on GPU or NPU accelerators.
\end{enumerate}

\subsection{Future Research Directions}
\label{sec:future_work}
Building upon these empirical insights, future work will focus on:
\begin{itemize}
    \item \textbf{Multi-Seed and Expanded Facility Evaluation:} Extending the benchmark across multiple random seeds and 50+ training epochs on expanded real-world warehouse facilities, incorporating targeted oversampling for sparse tail classes (\texttt{pallet\_truck}) to establish rigorous statistical confidence intervals ($\pm \sigma$).
    \item \textbf{Parallel Prefix Scans via TensorRT Plugins:} Implementing fused 2D selective scans using parallel associative prefix operators on NVIDIA Jetson Orin platforms, coupled with FP16 and INT8 quantization to bridge the gap between memory compliance and real-time execution.
    \item \textbf{Temporal State Propagation for Tracking:} Extending the 2D spatial scan across sequential video frames to propagate object state representations across time, improving multi-camera tracking stability in congested logistics corridors.
\end{itemize}

\section*{Acknowledgments}
The authors acknowledge the open-source computer vision and deep learning research communities, including the Ultralytics development team and NVIDIA Physical AI SmartSpaces contributors, for providing reproducible benchmarks and foundational detector architectures.
